\documentclass[11pt]{article}
\usepackage{amsmath,amsthm,amssymb}
\usepackage[margin=1in]{geometry}
\usepackage{booktabs,hyperref}
\emergencystretch=3em

\newtheorem{theorem}{Theorem}
\newtheorem{proposition}[theorem]{Proposition}
\newtheorem{corollary}[theorem]{Corollary}
\newtheorem{lemma}[theorem]{Lemma}
\theoremstyle{definition}
\newtheorem{definition}[theorem]{Definition}
\newtheorem{example}[theorem]{Example}
\newtheorem{remark}[theorem]{Remark}

\newcommand{\Hom}{\operatorname{Hom}}
\newcommand{\Ext}{\operatorname{Ext}}
\newcommand{\Aut}{\operatorname{Aut}}
\newcommand{\id}{\mathrm{id}}
\newcommand{\ob}{\operatorname{ob}}
\newcommand{\Ab}{\mathbf{Ab}}
\newcommand{\Sk}{\mathrm{Sk}}
\newcommand{\C}{\mathcal{C}}
\newcommand{\D}{\mathcal{D}}
\newcommand{\Z}{\mathbb{Z}}
\newcommand{\CatB}{\mathbf{Cat}_B}
\newcommand{\Igp}{\mathbf{I}}
\newcommand{\Bpi}{B\pi}
\newcommand{\infl}{\mathrm{inf}}
\newcommand{\res}{\mathrm{res}}
\newcommand{\tg}{\mathrm{tg}}
\newcommand{\im}{\operatorname{im}}
\newcommand{\loops}{\mathrm{loops}}

\title{Holonomy and prunability as the two edges of one spectral sequence\\
\large the Zappa--Sz\'ep composition holonomy is the transgression, on Ferri's $K_{2,3}$}
\author{MacBeth\\\small Kodamai}
\date{17 September 2026}

\begin{document}
\maketitle

\begin{center}\small
\fbox{\begin{minipage}{0.92\textwidth}\centering
\emph{Author:} MacBeth (Kodamai).\qquad\emph{Date:} 17 September 2026.\\[2pt]
\emph{Corresponds to} \texttt{scotmacbeth/work-in-progress} \emph{content-commit} \texttt{WIPHASH}.\\[2pt]
\textbf{Draft --- work in progress; Kodamai grant deliverable (Theory strand).}
For Rick (skew-brace / $H^2$ referee), then Neil. The structural results
(\S\ref{sec:reduce}--\S\ref{sec:unify}) are proved; the $K_{2,3}$ material of
\S\ref{sec:k23} is computed and machine-checked on the fibre edge, and the general
fibre-edge identification of prunability is flagged \textbf{open} throughout
(\S\ref{sec:status}).
\end{minipage}}
\end{center}

\begin{abstract}
A connected groupoid arising from a Zappa--Sz\'ep (matched-pair) factorisation carries
two proved degree-two obstruction theories: the multiplicative \emph{composition
holonomy} $[\omega]$, whose vanishing is exactly the existence of the global factorisation,
and Ferri's \emph{prunability} obstruction, whose vanishing lets the groupoid be pruned to
its loop bundle. Ferri asked (Remark~3.20 of his \emph{dynamical skew braces} paper) for a
group-theoretic or cohomological recast of prunability. We give one on his own example, and
locate both obstructions inside a single object. When the factorisation is normal, the categorical
cohomology of the groupoid is the abutment of the Lyndon--Hochschild--Serre spectral
sequence of the extension $1\to\D\to\pi\to\Gamma\to1$; its two degree-two edges are exactly
the two obstruction theories --- the base edge $H^2(\Gamma;M^\D)$ carries $[\omega]$, the
fibre edge $H^2(\D;M)^\Gamma$ carries prunability --- and the transgression $d_2(\varphi)
=\varphi_*[\omega]$ \emph{is} the holonomy. Consequently the two faces are never independent
summands: a nonzero base class is a differential, not a summand, so the naive direct sum
fails precisely when it would be interesting. On the fibre side we compute the fibre-edge
class on Ferri's own $K_{2,3}$: it is connected (so $[\omega]=0$) yet non-prunable, and the
fibre edge carries a nonzero class $[\pi]$, while a prunable Klein control kills both edges.
All finite claims are machine-checked; the general identification of prunability with the
fibre edge for arbitrary stars is stated as an open problem.
\end{abstract}

\section{Introduction}\label{sec:intro}

\paragraph{The problem.} A directed container, a matched pair of groups, and a Zappa--Sz\'ep
product are three names for the same move: build a large composite from two interacting
factors that act on one another. Whether the composite exists globally, and whether it can
be simplified, are the two questions a compositional-correctness theory must answer. For the
groupoid $G$ presented by a (quiver) skew brace factored over a vertex bundle, both questions
have been given cohomological answers, but by \emph{different} theories:
\begin{itemize}
\item[\textbf{(H)}] the \emph{composition holonomy} $[\omega]\in H^2(\Sk_\C;\D)$ over the
orbit base $\Sk_\C$, whose vanishing is equivalent to the existence of the global
Zappa--Sz\'ep factorisation $K=\C\bowtie\D$ \cite{gobstr,gzs}; and
\item[\textbf{(P)}] Ferri's \emph{prunability} obstruction \cite{ferri}, an additive
condition on the loop bundle (the loops over each object must form an additive subgroup),
whose vanishing lets $G$ be pruned to its loop subgroupoid, and which lives in the
direction-functor cohomology $H^2_{\CatB}(G)\cong H^2(\pi)$ of the vertex group
\cite{adm,sep2}.
\end{itemize}
These were proved \emph{separate}: $[\omega]$ vanishes identically for a full loop bundle,
so it cannot be prunability \cite{sep}. Ferri, working from the skew-brace side, asked for
the missing structural picture explicitly.

\begin{quote}\small
\textbf{Ferri, Remark 3.20 (paraphrased) \cite{ferri}.} \emph{Is there a group-theoretic or
cohomological description of the prunability obstruction, of the kind available for the
associated extension classes?}
\end{quote}

\paragraph{The gap.} ``Separate'' is not a structure. Knowing that $[\omega]$ and
prunability are two different classes leaves open the natural synthesis: are they two
independent summands of one categorical cohomology,
\begin{equation}\label{eq:conj}
   H^2_{\mathrm{cat}}(G;M)\;\overset{?}{\cong}\;
   \underbrace{H^2(\pi;M)}_{\text{prunability}}\;\oplus\;
   \underbrace{H^2_{\mathrm{base}}(\Sk_\C;M)}_{[\omega]},
   \qquad\text{both summands possibly nonzero?}
\end{equation}
This ``product decomposition'' is what one would write down first. It is false, and its
failure is the point.

\paragraph{The results.} We settle \eqref{eq:conj} in general and answer Ferri's question on
his own example, by putting both obstructions on one page. The structural statement below is
proved; the answer to Remark~3.20 is computed on $K_{2,3}$, and the general identification of
prunability with the fibre edge is left open (\S\ref{sec:status}).

\begin{theorem}[Structural, \S\ref{sec:reduce}--\S\ref{sec:unify}]\label{thm:main-struct}
Let $G$ be a connected groupoid whose vertex group $\pi$ has a normal abelian subgroup
$\D$ (the vertex bundle of a Zappa--Sz\'ep factorisation), with quotient $\Gamma=\pi/\D$
the vertex group of the orbit base. Then $H^2_{\mathrm{cat}}(G;A)\cong H^2(\pi;M)$ is the
abutment of the Lyndon--Hochschild--Serre spectral sequence
$E_2^{p,q}=H^p(\Gamma;H^q(\D;M))$, in which
\begin{enumerate}
\item[\textup{(i)}] the \emph{base edge} $E_2^{2,0}=H^2(\Gamma;M^\D)$ is the home of the
holonomy $[\omega]$;
\item[\textup{(ii)}] the \emph{fibre edge} $E_2^{0,2}=H^2(\D;M)^\Gamma$ is the home of
prunability;
\item[\textup{(iii)}] the transgression $d_2\colon E_2^{0,1}\to E_2^{2,0}$ \emph{is} the
holonomy: $d_2(\varphi)=\varphi_*[\omega]$.
\end{enumerate}
Consequently the direct sum \eqref{eq:conj} holds if and only if $[\omega]=0$; whenever
$[\omega]\neq0$ and some $\varphi$ has $\varphi_*[\omega]\neq0$, the base edge does not
survive and no such sum exists. The two obstructions are never simultaneously present as
independent summands.
\end{theorem}

The one-line reading of Theorem~\ref{thm:main-struct}: \emph{a nonzero holonomy is not a
summand next to prunability --- it is the differential that transgresses itself away.} This
is the group-theoretic recast Ferri asked for: prunability is the fibre edge, the holonomy
is the transgression, and both are edges of one extension's spectral sequence.

The theorem describes what happens on each edge in the abstract. To see that the fibre edge
is \emph{genuinely occupied} --- that prunability is a real, nonzero, computable class and
not a formal placeholder --- we compute it on Ferri's own example.

\begin{theorem}[Witness, computed, \S\ref{sec:k23}]\label{thm:main-witness}
On Ferri's $K_{2,3}$ \cite[Ex.~4.29]{ferri} --- a connected degree-two groupoid on
$\Z/4$ that is \emph{non-prunable} (Ferri's Counterexample~5.2) --- the two edges decouple:
the base edge vanishes, $[\omega]=0$ \textup{(}$G$ is connected, so the orbit base is
contractible\textup{)}, while the fibre edge carries a nonzero class
$[\pi]\in H^2(\Z/2;\Z/2)=\Z/2$. A prunable Klein-four control kills both edges. Thus the
fibre edge is the genuine detector of prunability, decoupled from the holonomy, on Ferri's
actual numbers.
\end{theorem}

\paragraph{Method.} Normality turns the whole question into ordinary group cohomology along
the extension $1\to\D\to\pi\to\Gamma\to1$ (\S\ref{sec:reduce}); the two obstruction theories
are then read off as the two axes of the LHS $E_2$-page (\S\ref{sec:lhs}); a one-page cocycle
computation identifies the transgression with the extension class (\S\ref{sec:tg}); and a
survival count gives the collapse locus and the no-go (\S\ref{sec:nogo}). The $K_{2,3}$
computation (\S\ref{sec:k23}) is the additive twin of the multiplicative Schreier line for
$[\omega]$: the fibre-edge class is the nonsplitting of the additive extension
$0\to\{0,2\}\to\Z/4\to\Z/2\to0$ read at each star, and it is machine-checked in Lean.

\paragraph{Consequences.} Three previously separate results --- the holonomy face (H), the
prunability face (P), and their proved separation (S) \cite{sep,sep2} --- become one object.
The separation theorem is recovered as the contractible-base slice $\Gamma\simeq\ast$, where
the base edge is forced to zero and only the fibre edge survives; $K_{2,3}$ is exactly this
slice, realised on a real example. For the grant Theory strand this is the honest
``internal replacement / Zappa--Sz\'ep, fully cohomological'' statement: compositional
correctness of the directed-container composite is $[\omega]=0$, which is the collapse of one
spectral sequence whose two edges are the base handoff geometry and the additive prunability.

\paragraph{Outline.} \S\ref{sec:prelim} fixes the setup and defines $[\omega]$ and
prunability. \S\ref{sec:reduce} reduces to a group extension. \S\ref{sec:lhs} identifies the
two edges. \S\ref{sec:tg} proves the transgression is the holonomy. \S\ref{sec:nogo} gives
the collapse locus, the no-go, and the finite witnesses. \S\ref{sec:k23} computes the fibre
edge on $K_{2,3}$ and answers Remark~3.20. \S\ref{sec:unify} assembles the unification.
\S\ref{sec:status} states the honest scope, including the open general identification.

\section{Preliminaries}\label{sec:prelim}

We keep only what the proofs use. Throughout, $G$ is a connected groupoid with a chosen base
object $x_0$ and vertex group $\pi=G(x_0,x_0)$; $A$ is a $G$-module and $M:=A(x_0)$ its
value at $x_0$. All cohomology is group cohomology unless decorated. The
Lyndon--Hochschild--Serre (LHS) spectral sequence of a group extension and its five-term
exact sequence are used as in Weibel \cite[\S6.8]{weibel}; the transgression formula in the
general twisted case is Neukirch--Schmidt--Wingberg \cite{nsw}.

\paragraph{Zappa--Sz\'ep factorisation and the holonomy $[\omega]$.} A Zappa--Sz\'ep (matched
pair) factorisation presents the composite groupoid $K$ as an interaction $\C\bowtie\D$ of a
base groupoid $\C$ and a vertex bundle $\D$ that act on each other. The obstruction to the
global factorisation existing is a degree-two class, the \emph{composition holonomy}
$[\omega]\in H^2(\Sk_\C;\D)$ over the orbit base $\Sk_\C$ (the skeleton of $\C$), and
$K=\C\bowtie\D$ exists iff $[\omega]=0$ \cite{gobstr,gzs}. When $\D$ is normal in $\pi$ this
is the Schreier class of the extension $1\to\D\to\pi\to\Gamma\to1$ with $\Gamma=\Sk_\C$, and
it is genuinely nonzero --- e.g.\ $[\omega]\neq0$ for $\Z/4$ over $\Z/2$ \cite{gzs}. When the
loop bundle is \emph{full}, the orbit base is connected, hence contractible, and
$[\omega]=0$ identically \cite{sep}.

\paragraph{Prunability.} For a skew-brace groupoid, the loops over an object $x$ form a set
$\loops(x)$ inside the additive group at $x$. Ferri calls $G$ \emph{prunable} when the loop
subgroupoid is a $\rightharpoonup$-ideal --- concretely, when $\loops(x)$ is an additive
subgroup for every $x$ --- so that $G$ may be pruned to its loop bundle \cite{ferri}. The
obstruction is additive and lives on the vertex group: the direction-functor cohomology of
Ambra--Duvieusart--Montoli \cite{adm} reduces, for a connected groupoid, to
$H^2_{\CatB}(G;A)\cong H^2(\pi;M)$ \cite[T1]{sep2}, the additive face that carries
prunability \cite{sep}. Ferri's Remark~3.20 asks for the cohomological picture behind this
condition; \S\ref{sec:k23} supplies it on his own example.

\paragraph{The separation, and the synthesis it invites.} The two faces were proved
orthogonal: $[\omega]\perp$ prunability, from a computed witness \cite{sep} and structurally
from the decomposition $G\cong\Igp(\ob G)\times\Bpi$ in the full-bundle case \cite[T2]{sep2}.
Orthogonality is what makes the direct sum \eqref{eq:conj} tempting. But \eqref{eq:conj}
silently uses one symbol $\Sk_\C$ for two different bases:
\begin{itemize}
\item In the \emph{connected/full-loop-bundle} regime, $\Sk_\C=\Igp(\ob G)$ is the coarse
groupoid on the objects, a disjoint union of contractible components, so
$H^{\ge1}(\Sk_\C;-)=0$ always and the base summand of \eqref{eq:conj} is identically zero.
Here $G\cong\Igp(\ob G)\times\Bpi$ is a genuine product and \eqref{eq:conj} says nothing but
$H^2_{\mathrm{cat}}=H^2(\pi)$.
\item In the \emph{proper-normal-bundle} regime, $\Sk_\C=B\Gamma$ is the classifying groupoid
of a quotient group and is generally not contractible; this is where $[\omega]\neq0$. But now
$G$ is \emph{not} a product: $(\D,\Gamma,\pi)$ form a non-split \emph{extension}.
\end{itemize}
So the two regimes never coexist inside one product. The honest object that sees a
non-contractible base is a group extension, and its cohomology is governed by LHS, not by
K\"unneth. We now show $[\omega]$ appears there as the transgression.

\section{Reduction to a group extension}\label{sec:reduce}

We work in the normal Zappa--Sz\'ep regime of \cite{gzs}, the only regime in which the base
carries a nonzero class.

\begin{proposition}[Reduction, {\cite[Prop.~2.1--2.2]{gzs}}, {\cite[T1]{sep2}}]
\label{prop:reduce}
Let $G$ be a connected groupoid whose vertex group $\pi$ admits a normal abelian subgroup
$\D\trianglelefteq\pi$ (the vertex bundle of a Zappa--Sz\'ep factorisation), with quotient
$\Gamma=\pi/\D$ the vertex group of the orbit base $\Sk_\C=B\Gamma$. Then:
\begin{enumerate}
\item[\textup{(a)}] $\D$ is a $\Gamma$-module via conjugation, and the Zappa--Sz\'ep holonomy
$[\omega]\in H^2(\Sk_\C;\D)=H^2(\Gamma;\D)$ is the Schreier class of the extension
\begin{equation}\label{eq:ext}
   1\longrightarrow\D\longrightarrow\pi\xrightarrow{\ q\ }\Gamma\longrightarrow1;
\end{equation}
\item[\textup{(b)}] for a $G$-module $A$ with $M:=A(x_0)$, the categorical
(direction-functor) cohomology reduces to the vertex group,
$H^2_{\mathrm{cat}}(G;A)=H^2_{\CatB}(G;A)\cong H^2(\pi;M)$.
\end{enumerate}
\end{proposition}
\begin{proof}
(a) is \cite[Prop.~2.1--2.2]{gzs}: normality gives $\Sk_\C=\pi/\D=\Gamma$ and the defect
$2$-cochain is the Schreier factor set of \eqref{eq:ext}. (b) is \cite[T1]{sep2},
$H^n_{\CatB}(G,A)\cong H^n(\pi,M)$, whose $n=2$ instance rests on the reconstructed Schreier
correspondence there and elementary homological algebra.
\end{proof}

Proposition~\ref{prop:reduce} moves the whole question into ordinary group cohomology: how
does $H^2(\pi;M)$ decompose along \eqref{eq:ext}, and where do the two faces sit? The answer
is the LHS spectral sequence.

\section{The spectral sequence and its two edges}\label{sec:lhs}

\begin{theorem}[LHS decomposition of the categorical cohomology]\label{thm:lhs}
With the data of Proposition~\ref{prop:reduce}, $H^\bullet_{\mathrm{cat}}(G;A)\cong
H^\bullet(\pi;M)$ is the abutment of the first-quadrant cohomological spectral sequence
\begin{equation}\label{eq:E2}
   E_2^{p,q}\;=\;H^p\bigl(\Gamma;\,H^q(\D;M)\bigr)\;\Longrightarrow\;H^{p+q}(\pi;M),
\end{equation}
of the extension \eqref{eq:ext} \cite[\S6.8]{weibel}. Its two degree-two edge terms are
exactly the two obstruction faces:
\[
   \begin{aligned}
   E_2^{2,0}&=H^2\bigl(\Gamma;M^\D\bigr)
   &&\text{(\emph{base edge}: the home of $[\omega]$)},\\
   E_2^{0,2}&=H^2(\D;M)^\Gamma
   &&\text{(\emph{fibre edge}: the prunability face)} .
   \end{aligned}
\]
The abutment is filtered with associated graded
$\mathrm{gr}\,H^2(\pi;M)=E_\infty^{2,0}\oplus E_\infty^{1,1}\oplus E_\infty^{0,2}$,
where $E_\infty^{p,q}$ is the subquotient of $E_2^{p,q}$ surviving the differentials.
\end{theorem}
\begin{proof}
The LHS spectral sequence of a group extension with abelian coefficient module $M$ is
standard \cite[Thm.~6.8.2]{weibel}; \eqref{eq:E2} is its $E_2$-page. The edge
identifications $E_2^{p,0}=H^p(\Gamma;M^\D)$ and $E_2^{0,q}=H^q(\D;M)^\Gamma$ are the
definition of the two axes. That the base edge is where $[\omega]$ lives is
Proposition~\ref{prop:reduce}(a) ($[\omega]\in H^2(\Gamma;\D)$, mapped to $H^2(\Gamma;M^\D)$
by any $\Gamma$-map $\D\to M^\D$); that the fibre edge is the vertex/prunability face is
\cite{sep,sep2}. The three-step filtration in total degree $2$ is the standard convergence
of a first-quadrant sequence.
\end{proof}

The two faces are not \emph{a priori} summands: they are two edges of one page, and whether
they survive to the abutment is decided by the differentials. The relevant one is the
transgression, and it is the holonomy itself.

\section{The transgression is the holonomy}\label{sec:tg}

\begin{lemma}[The transgression is cup with the holonomy]\label{lem:tg}
In the five-term exact sequence of \eqref{eq:E2},
\[
   0\to H^1(\Gamma;M^\D)\xrightarrow{\infl}H^1(\pi;M)\xrightarrow{\res}
   H^1(\D;M)^\Gamma\xrightarrow{\ \tg\ }H^2(\Gamma;M^\D)\xrightarrow{\infl}H^2(\pi;M),
\]
the transgression $\tg=d_2^{0,1}$ is given by pushforward of the extension class: for a
$\Gamma$-equivariant homomorphism $\varphi\in\Hom_\Gamma(\D,M)=H^1(\D;M)^\Gamma$,
\[
   \tg(\varphi)\;=\;\varphi_*[\omega]\;\in\;H^2(\Gamma;M^\D),
\]
where $\varphi_*\colon H^2(\Gamma;\D)\to H^2(\Gamma;M^\D)$ is induced by $\varphi$ and
$[\omega]$ is the Schreier class of \eqref{eq:ext}.
\end{lemma}
\begin{proof}
We give the cocycle computation for the case $M$ a trivial $\pi$-module (which covers all
witnesses below and the $\Z/4$-Schreier line); the general $\Gamma$-module case is
\cite[transgression formula]{nsw}, identical after keeping the covariant end. Choose a
normalised set section $s\colon\Gamma\to\pi$ of $q$, $s(1)=1$, with Schreier factor set
\[
   s(a)\,s(b)=\omega(a,b)\,s(ab),\qquad \omega(a,b)\in\D,\quad [\omega]\in H^2(\Gamma;\D).
\]
An element of $E_2^{0,1}=H^1(\D;M)^\Gamma$ is a $\Gamma$-invariant homomorphism
$\varphi\colon\D\to M$. Running the construction of the transgression in the double complex
of the extension \cite[\S6.8]{weibel}: $\varphi$ lifts to the $1$-cochain
$\tilde\varphi\in C^1(\pi;M)$, $\tilde\varphi(g)=\varphi(\text{$\D$-part of }g)$ along the
section, and its image under the vertical differential, pushed to the base, is the
$2$-cochain
\[
   (a,b)\longmapsto \varphi\bigl(\omega(a,b)\bigr)=(\varphi_*\omega)(a,b).
\]
Thus $\tg(\varphi)=[\varphi_*\omega]=\varphi_*[\omega]$. (Concretely: $\tg(\varphi)$ is the
obstruction to extending $\varphi$ from $\D$ to a homomorphism $\pi\to M$; it is $0$ iff the
pushed-out extension $1\to M\to\varphi_*\pi\to\Gamma\to1$ splits, i.e.\ iff
$\varphi_*[\omega]=0$.)
\end{proof}

\section{Collapse and the no-go}\label{sec:nogo}

\begin{corollary}[When the clean sum holds]\label{cor:collapse}
The naive direct sum \eqref{eq:conj} of the two edges equals $H^2(\pi;M)$ whenever
\eqref{eq:E2} collapses at $E_2$ in total degree $2$ with vanishing cross term $E_2^{1,1}$.
Each of the following is sufficient:
\begin{enumerate}
\item[\textup{(i)}] \textbf{Coarse/connected base} $\Gamma=1$: then $E_2^{p,q}=0$ for $p>0$,
the base edge and cross term vanish, and $H^2(\pi;M)=E_2^{0,2}=H^2(\D;M)$. Since here
$\D=\pi$ this is the identity $H^2_{\mathrm{cat}}(G)=H^2(\pi)$ --- the clean sum degenerates
to its fibre summand, and $[\omega]=0$.
\item[\textup{(ii)}] \textbf{Direct product} $\pi=\D\times\Gamma$: then $[\omega]=0$ and the
sequence collapses (K\"unneth), giving the \emph{full} associated graded $H^2(\pi;M)\cong
E_2^{2,0}\oplus E_2^{1,1}\oplus E_2^{0,2}$ --- the honest K\"unneth form, including the cross
term $H^1(\Gamma)\otimes H^1(\D)$ that \eqref{eq:conj} omits.
\item[\textup{(iii)}] \textbf{Coprime orders} $\gcd(|\D|,|\Gamma|)=1$ (Schur--Zassenhaus):
$E_2^{p,q}=0$ for $p,q>0$, so $H^2(\pi;M)=E_2^{2,0}\oplus E_2^{0,2}$ and $[\omega]=0$.
\end{enumerate}
In every collapsing case the base class is forced to zero: $[\omega]=0$.
\end{corollary}
\begin{proof}
(i) A coarse groupoid is equivalent to the terminal category, so $H^{\ge1}(\Gamma;-)=0$; the
surviving column is $p=0$. (ii) A direct product splits \eqref{eq:ext}, so $[\omega]=0$;
collapse is the Eilenberg--Zilber/K\"unneth theorem for $B\D\times B\Gamma$. (iii) For coprime
orders the higher cohomology of each factor is annihilated by both $|\D|$ and $|\Gamma|$,
hence $0$ off the edges; a split extension (Schur--Zassenhaus) has $[\omega]=0$.
\end{proof}

Corollary~\ref{cor:collapse} is the precise content of \eqref{eq:conj}: it is \emph{true
exactly on the locus $[\omega]=0$}, and there it is either the trivial statement
$H^2_{\mathrm{cat}}=H^2(\pi)$ (connected case, the prunability-only corner realised by
$K_{2,3}$ in \S\ref{sec:k23}) or the genuine K\"unneth sum (direct product) --- but in
neither case is the base summand nonzero. The interesting ``both nonzero'' regime is empty.

\begin{theorem}[{No clean sum when $[\omega]\neq0$}]\label{thm:nogo}
Suppose $[\omega]\neq0$ and there is a $\Gamma$-equivariant $\varphi\colon\D\to M$ with
$\varphi_*[\omega]\neq0$ in $H^2(\Gamma;M^\D)$ \textup{(}e.g.\ $M=\D$, $\varphi=\id$,
whenever $[\omega]\neq0$\textup{)}. Then:
\begin{enumerate}
\item[\textup{(a)}] the transgression $d_2^{0,1}=\tg$ is nonzero
(Lemma~\ref{lem:tg}), so the base edge does not survive: $E_\infty^{2,0}\subsetneq
E_2^{2,0}$, and the image $\varphi_*[\omega]$ of the holonomy is \emph{killed} in the
abutment;
\item[\textup{(b)}] consequently $H^2(\pi;M)$ is \emph{not} the direct sum \eqref{eq:conj} of
the two edges: $\dim H^2(\pi;M)<\dim E_2^{2,0}+\dim E_2^{1,1}+\dim E_2^{0,2}$ whenever
$\tg\neq0$.
\end{enumerate}
In particular the two faces are \emph{never} simultaneously present as independent summands
of $H^2_{\mathrm{cat}}(G)$: a nonzero base class is not a summand but a differential.
\end{theorem}
\begin{proof}
(a) By Lemma~\ref{lem:tg}, $\tg(\varphi)=\varphi_*[\omega]\neq0$, so $d_2^{0,1}\neq0$. The
differential $d_2^{0,1}\colon E_2^{0,1}\to E_2^{2,0}$ has nonzero image, hence
$E_3^{2,0}=E_2^{2,0}/\im d_2^{0,1}$ is a proper quotient; as the base edge receives no
further differentials ($d_r^{p,0}$ into $(2,0)$ come from $(2-r,r-1)$, a negative column for
$r\ge3$), $E_\infty^{2,0}=E_3^{2,0}\subsetneq E_2^{2,0}$, and the class $\varphi_*[\omega]$
(a generator of $\im d_2$) is annihilated. (b) The abutment has
$\dim H^2=\sum_{p+q=2}\dim E_\infty^{p,q}\le\sum_{p+q=2}\dim E_2^{p,q}$ with strict
inequality once any $E_\infty^{p,q}\subsetneq E_2^{p,q}$; part (a) gives this at $(2,0)$.
\end{proof}

\begin{corollary}[Finite witnesses, machine-checked]\label{cor:witness}
For trivial $\Z/2$-coefficients $M=\mathbb F_2$
(\texttt{2026-09-17-lhs-}\allowbreak\texttt{decomposition-check.py}):
\begin{center}\footnotesize\setlength{\tabcolsep}{4pt}
\begin{tabular}{lcccccl}
\toprule
$\pi=\D.\Gamma$ & split? & $[\omega]$ & $\dim H^2(\pi)$ &
$\dim(E_2^{2,0}\!\oplus\!E_2^{1,1}\!\oplus\!E_2^{0,2})$ & $\mathrm{rk}\,\tg$ & verdict\\
\midrule
$\Z/4=\Z/2.\Z/2$          & no  & $\neq0$ & $1$ & $1{+}1{+}1=3$ & $1$ & \textbf{no sum} (base edge dies)\\
$Q_8=(\Z/2)^2.\Z/2$       & no  & $\neq0$ & $2$ & $3{+}2{+}1=6$ & $1$ & \textbf{no sum}\\
$\Z/2\times\Z/2=\Z/2.\Z/2$& yes & $0$     & $3$ & $1{+}1{+}1=3$ & $0$ & clean (full K\"unneth)\\
$\Z/6=\Z/2.\Z/3$          & yes & $0$     & $1$ & $0{+}0{+}1=1$ & $0$ & clean (fibre only)\\
\bottomrule
\end{tabular}
\end{center}
For $\Z/4$ the transgression is an isomorphism $E_2^{0,1}\xrightarrow{\sim}E_2^{2,0}$, so the
entire base edge $H^2(\Z/2;\mathbb F_2)$ --- the home of the nonzero Schreier $[\omega]$ of
$\Z/4$ over $\Z/2$ \cite{gzs} --- \emph{vanishes} from the abutment. The naive sum overcounts
by $2$; the true $H^2(\Z/4;\mathbb F_2)$ is $1$-dimensional.
\end{corollary}

\section{The fibre edge is occupied: Ferri's \texorpdfstring{$K_{2,3}$}{K23}}\label{sec:k23}

Theorem~\ref{thm:nogo} shows the base edge cannot be an independent summand. It says nothing
about whether the fibre edge is ever nonzero on a real skew brace --- whether prunability is
an actual class or a formal slot. We settle this on Ferri's own example, which is
\emph{connected} (so we are in the collapse corner \ref{cor:collapse}(i), $[\omega]=0$) yet
\emph{non-prunable}. This is the honest place to test prunability: it is the corner where the
holonomy is forced to zero, so any obstruction that remains lives on the fibre edge alone.

\paragraph{The example.} Ferri's $K_{2,3}$ \cite[Ex.~4.29, Table~3]{ferri} has underlying set
$A=\Z/4$ and two objects $S_2,S_3$; the direction data give a connected degree-two groupoid.
Its loop bundle is
\[
   \loops(S_2)=\{0,3\},\qquad \loops(S_3)=\{0,1\}.
\]
Neither is an additive subgroup of $(\Z/4,+)$ ($3+3=2\notin\{0,3\}$, $1+1=2\notin\{0,1\}$),
so $G$ is \emph{not} prunable --- this is Ferri's Counterexample~5.2.

\begin{remark}[The quasigroup subtlety, stated plainly]\label{rmk:quasigroup}
The star operations $\bullet_{S_2},\bullet_{S_3}$ of Table~3 are left \emph{quasigroups}, not
groups: their rows are permutations but their columns are not, and both are non-commutative
and non-associative (machine-checked: e.g.\ $\bullet_{S_2}$ has $(1\bullet1)\bullet1=1\neq
3=1\bullet(1\bullet1)$). Since every group of order $4$ is abelian, $\bullet_\lambda$ cannot
itself be the additive group. The skew brace's genuine additive group on $A$ is the ordinary
$(\Z/4,+)$ (Ferri's convention $a\cdot b=a+b$), and it is with respect to $(\Z/4,+)$ that
``additive subgroup'', hence prunability, is meaningful. The loop sets \emph{are} closed
under the quasigroup $\bullet_\lambda$ (they are sub-left-quasigroups), but that is not the
condition Ferri's prunability tests; the defect is precisely that they are not subgroups of
$(\Z/4,+)$. We flag this because the identification of the fibre-edge class below with
prunability rests on this $(\Z/4,+)$ reading (see \S\ref{sec:status}).
\end{remark}

\paragraph{Base edge $E_2^{2,0}=H^2(\Gamma;M^\D)$.} $K_{2,3}$ is connected, so the orbit
base is the coarse groupoid on $\{S_2,S_3\}$, which is contractible: $H^{\ge1}(\Gamma;-)=0$
for any coefficients. Hence $[\omega]=0$. This is rigorous and coefficient-independent
(Corollary~\ref{cor:collapse}(i) on this example).

\paragraph{Fibre edge $E_2^{0,2}=H^2(\D;M)^\Gamma$.} With $\Gamma\simeq\ast$ the $\Gamma$-in\-vari\-ants
are everything, and the fibre edge is $H^2(\D;M)$ computed at each star. Read as the additive
Schreier class of the extension
\[
   0\longrightarrow A_2\longrightarrow(\Z/4,+)\longrightarrow \Z/2\longrightarrow0,
   \qquad A_2=\{0,2\}\cong\Z/2,
\]
we have $H^2(\Z/2;\Z/2)=\Z/2$, and the extension is \emph{nonsplit} at both $S_2$ and $S_3$:
its Schreier cocycle is the carry $f(1,1)=2\in A_2$, which is not a coboundary. Hence the
fibre-edge class
\[
   [\pi]_{S_2}=[\pi]_{S_3}=\text{the nonzero generator of }H^2(\Z/2;\Z/2)=\Z/2 .
\]
This computation is machine-checked in Lean, sorry-free
(\texttt{2026-09-17-}\allowbreak\texttt{fibre-h2-z2.lean}): $H^2(\Z/2;\Z/2)=\Z/2$ as the cokernel of the norm;
the $\Z/4$ carry cocycle $f(a,b)=a\wedge b$ is a genuine cocycle and \emph{not} a coboundary
(matching at $(0,0)$ forces $g(0)=0$, whence $\delta^1 g(1,1)=0\neq1$); and the Klein control
cocycle is $0$.

\paragraph{The control.} Replacing the $\Z/4$ star by the Klein four-group $V_4$ keeps the
groupoid connected but makes the loops an additive subgroup, so $G$ becomes prunable. Now the
additive extension $0\to\Z/2\to V_4\to\Z/2\to0$ \emph{splits}, so the fibre-edge class is
zero. Both edges are dead. The control isolates the fibre edge as the genuine detector of
prunability: it is nonzero exactly when $G$ fails to be prunable.

\begin{center}\small
\begin{tabular}{lccccc}
\toprule
case & prunable? & base $H^2$ & $[\omega]$ & star extension & $[\pi]$ (fibre edge)\\
\midrule
Ferri $K_{2,3}$ ($\Z/4$ stars) & \textbf{no}  & $0$ & $0$ & nonsplit (both $S_2,S_3$) & $\neq0$ (both)\\
Klein control ($V_4$ star)     & yes          & $0$ & $0$ & split                      & $0$\\
\bottomrule
\end{tabular}
\end{center}

\paragraph{Answer to Remark~3.20.} On Ferri's own $K_{2,3}$, the prunability obstruction
\emph{is} the fibre-edge class $[\pi]\in H^2(\D;M)^\Gamma$ of the LHS spectral sequence
\eqref{eq:E2}, equivalently the nonsplitting of the additive star extension
$0\to A_2\to(\Z/4,+)\to\Z/2\to0$. This is the additive twin of the multiplicative Schreier
line that carries $[\omega]$: $[\omega]$ is the nonsplitting of the multiplicative extension
$1\to\D\to\pi\to\Gamma\to1$ on the base edge, and prunability is the nonsplitting of the
additive extension on the fibre edge. The dividing line, $\Z/4$ (nonsplit) versus Klein
(split), is exactly the additive twin of the $\Z/4$-versus-$\Z/2\times\Z/2$ line for
$[\omega]$ in Corollary~\ref{cor:witness}. This is the group-theoretic/cohomological
description Ferri asked for.

\section{What the cluster is, correctly}\label{sec:unify}

The three cluster results unify into one object --- a spectral sequence, not a direct sum.

\begin{theorem}[Structural unification]\label{thm:unify}
Under Proposition~\ref{prop:reduce}, the categorical cohomology $H^2_{\mathrm{cat}}(G;A)\cong
H^2(\pi;M)$ is the total degree-two term of the LHS spectral sequence \eqref{eq:E2}, in which:
\begin{enumerate}
\item the \emph{fibre edge} $E_2^{0,2}=H^2(\D;M)^\Gamma$ is the additive/prunability face
\textup{(P)}, occupied by a nonzero class on Ferri's $K_{2,3}$ (\S\ref{sec:k23});
\item the \emph{base edge} $E_2^{2,0}=H^2(\Gamma;M^\D)$ is the holonomy face, receiving the
transgression $d_2$;
\item the \emph{transgression} $d_2\colon E_2^{0,1}\to E_2^{2,0}$ \emph{is} the Zappa--Sz\'ep
holonomy, $d_2(\varphi)=\varphi_*[\omega]$ \textup{(H)}.
\end{enumerate}
The proved separation $[\omega]\perp$ prunability \textup{(S)} is the special case
$\Gamma\simeq\ast$ (base contractible, only the fibre edge survives) --- precisely the corner
$K_{2,3}$ inhabits. In general the two faces are not orthogonal summands but the two axes of
one page, \emph{coupled} by the transgression, which is $[\omega]$ itself. The clean direct
sum \eqref{eq:conj} holds iff this coupling vanishes, i.e.\ iff $[\omega]=0$.
\end{theorem}
\begin{proof}
Assemble Theorem~\ref{thm:lhs} (edges), Lemma~\ref{lem:tg} (transgression $=$ holonomy),
Corollary~\ref{cor:collapse} (collapse $\Leftrightarrow$ the true locus of \eqref{eq:conj}),
Theorem~\ref{thm:nogo} (no-go off that locus), and \S\ref{sec:k23} (the fibre edge is
occupied). The separation is Corollary~\ref{cor:collapse}(i): $\Gamma=1\Rightarrow$ base edge
and cross term vanish, $H^2_{\mathrm{cat}}=H^2(\D)=H^2(\pi)$, orthogonality is the triviality
of the base axis.
\end{proof}

\begin{remark}[Why this is the right correction]
The conjecture \eqref{eq:conj} failed for a structural reason worth stating plainly: a
nonzero base class $[\omega]$ \emph{is} a non-split extension, and a non-split extension is
exactly a nonzero transgression. One cannot ``add'' the holonomy as a summand next to the
prunability class, because the holonomy is the mechanism by which the base cohomology is
glued into the fibre cohomology. The separation theorem saw the two faces as orthogonal only
because it lived on the contractible-base slice where the transgression is forced to zero;
off that slice, orthogonality is replaced by a single coupling differential. This is strictly
more informative, and it is the honest ``internal replacement / Zappa--Sz\'ep'' deliverable:
compositional correctness of the directed-container composite $=[\omega]=0=$ the collapse of
one spectral sequence, whose two edges are the base handoff geometry and the vertex/additive
prunability.
\end{remark}

\section{Status and honest scope}\label{sec:status}
\paragraph{Proved.} Theorem~\ref{thm:lhs} (LHS decomposition; standard SS $+$ proved edge
identifications), Lemma~\ref{lem:tg} (transgression $=$ holonomy; self-contained cocycle
proof for trivial coefficients, cited \cite{nsw} in general), Corollary~\ref{cor:collapse}
(collapse locus $=\{[\omega]=0\}$), Theorem~\ref{thm:nogo} (no-go), Theorem~\ref{thm:unify}
(unification). These rest on the proved Proposition~\ref{prop:reduce}(a)(b)
$=$ \cite{gzs,sep2}, the classical LHS machinery \cite{weibel,nsw}, and Schur--Zassenhaus.
\paragraph{Computed.} The four finite witnesses of Corollary~\ref{cor:witness}
(machine-checked bar complex over $\mathbb F_2$ and five-term rank count) and the $K_{2,3}$
decoupling table of \S\ref{sec:k23}. The fibre-edge kernel ($H^2(\Z/2;\Z/2)=\Z/2$, $\Z/4$
carry cocycle $\neq0=$ Klein) is machine-checked in Lean, sorry-free
(\texttt{2026-09-17-}\allowbreak\texttt{fibre-h2-z2.lean}).
\paragraph{Not claimed.}
\begin{itemize}
\item The clean direct sum \eqref{eq:conj} with both summands nonzero --- \textbf{refuted}
(Theorem~\ref{thm:nogo}).
\item \textbf{Open (the load-bearing caveat).} The general structural identification of
Ferri prunability with a specific element of the fibre edge $E_2^{0,2}$, for \emph{arbitrary}
stars and twisted coefficients, is \emph{not} proved here. We verify it on $K_{2,3}$ and its
Klein control, and on the two-object $\Z/2$ family generally; a structural proof that the
categorical LHS fibre edge literally receives Ferri prunability for a genuinely dynamical
(quasigroup-only) star remains open. On $K_{2,3}$ the identification additionally rests on
reading $+_\lambda$ as $(\Z/4,+)$ (Ferri's Counterexample~5.2 convention;
Remark~\ref{rmk:quasigroup}), since the star $\bullet_\lambda$ is quasigroup-only. This is the
separately tracked node \texttt{prunability-is-fibre-edge-class}.
\item The nonabelian / non-normal Zappa--Sz\'ep matched-pair regime (Kac/Pirashvili $H^2$) is
out of scope, as in \cite{gzs}.
\end{itemize}

\begin{thebibliography}{99}
\bibitem{gobstr} MacBeth, \emph{The global-closure obstruction (G) is a cohomology class},
Kodamai note, 11 June 2026.
\bibitem{gzs} MacBeth, \emph{Invertibility does not kill the Zappa--Sz\'ep merge obstruction:
the groupoid case is group cohomology}, Kodamai note, 24 July 2026
(\texttt{2026-07-24-groupoid-zs-obstruction.tex}).
\bibitem{sep} MacBeth, \emph{Ferri prunability is not the Zappa--Sz\'ep composition holonomy:
a separation theorem}, Kodamai note, 16 September 2026
(\texttt{2026-09-16-ferri-prunability-vs-zs-holonomy.tex}).
\bibitem{sep2} MacBeth, \emph{The direction-functor $H^2$ of a connected groupoid is its
vertex-group cohomology}, Kodamai note, 16 September 2026
(\texttt{2026-09-16-h2catb-vertex-group-reduction.tex}).
\bibitem{ferri} D.~Ferri, \emph{On dynamical skew braces and skew bracoids}, J.\ Pure Appl.\
Algebra \textbf{229} (2025), no.~9, art.~108036, arXiv:2410.10717
(Example~4.29 / Table~3, Counterexample~5.2, Remark~3.20).
\bibitem{adm} S.~Ambra, A.~Duvieusart, A.~Montoli, \emph{A direction functor approach to the
cohomology of small categories}, arXiv:2608.07380 (2026).
\bibitem{weibel} C.~A.~Weibel, \emph{An Introduction to Homological Algebra}, Cambridge Univ.\
Press, 1994 (\S6.8, the Lyndon--Hochschild--Serre spectral sequence and the five-term exact
sequence).
\bibitem{nsw} J.~Neukirch, A.~Schmidt, K.~Wingberg, \emph{Cohomology of Number Fields}, 2nd
ed., Springer, 2008 (Ch.~I, transgression and the extension class).
\end{thebibliography}
\end{document}
